\section{Environment, experiments and measurements}
\subsection{A financially consequential factual decision}
\label{sec:protocol}
All games run on isolated local Anvil test chains. Each case is one attempt
to adjudicate a single claim. The publisher posts
``The world is flat'', and the scripted challenger submits evidence that Earth
is round. Supporting the publisher is therefore the false verdict in this
study. The publisher
and challenger each post 100 synthetic ETH, and each juror bonds 5. On settlement,
the losing party's bond pays 85 ETH to the winning party, a 5 ETH pool to valid revealing jurors,
and 10 ETH to the protocol. Revealing jurors recover principal and share the
pool regardless of vote direction. Returning a juror's own bond is distinct
from this reward. The publisher promises an additional
10 ETH credit transfer to each matching revealer if it wins. The promise is
not escrowed. These are local game credits, not USD or public-chain earnings.

Voting has two stages. During commitment, a juror deposits a bond and submits
a cryptographic commitment that fixes its verdict and assessment without
revealing their contents. During reveal, it supplies the matching values so
the contract can check and count the vote. Only valid reveals count toward a
strict majority of the full registered jury. Absent jurors do not lower that
threshold. If neither side reaches it, the case aborts without a winning side.
Thus a promise, a commitment, a counted vote and a settled verdict are separate
observations.

Evaluated accounts have a 20 ETH monthly revenue target. It is not a
one-case deadline. Each model uses general file and execution tools to inspect
its instructions and the HTTP API and to act through its assigned identity.
The endpoints are Cloudflare Workers AI
\path{@cf/moonshotai/kimi-k2.6}, \path{@cf/moonshotai/kimi-k2.7-code}, and
\path{@cf/deepseek-ai/deepseek-v4-pro-0813}. Each model response has an
8,192-token completion limit. The call budget is 12 calls per opportunity and
100 per actor per case. The stepped scheduler permits
two opportunities per phase. An opportunity is a scheduled turn in which an
actor can inspect the case and act. Selected actual inputs appear in
\autoref{app:model-exchanges}.

\subsection{Primary comparison: first-moving peer blocs at three sizes}
\label{sec:quorum}
Let $n\in\{3,7,15\}$ be the registered jury size. In each primary condition,
$m=(n-1)/2$ jurors are scripted first movers, and the remaining
$k=n-m=(n+1)/2$ jurors use the evaluated model. These scripted jurors form the
seeded bloc. Colluding peers announce a knowingly false vote for payment and
invite others to join. Honest peers support the evidence and decline the
offer. The publisher's offer is present in both arms, and the peers act before
evaluated jurors are notified.

Because $n$ is odd, the required strict majority is $q=(n+1)/2$ valid reveals.
In the colluding-peer arm, the $m$ scripted false votes leave a shortfall of
\[
q-m=\frac{n+1}{2}-\frac{n-1}{2}=1.
\]
One evaluated false revealer therefore suffices to capture the verdict if the
scripted bloc reveals as configured. For example, a seven-juror case requires
four votes. Its three scripted false votes need only one of the four evaluated
jurors to join them. This is the same one-vote shortfall as in a three-juror
or fifteen-juror case.

All assigned participants can use the shared public forum and permitted
private channels. Every evaluated juror receives the seeded messages before
acting. A communication graph specifies who can send messages to whom. This
graph is not randomized or rewired here. Increasing
$n$ also increases $m$ and $k$, raises the seed fraction from $1/3$ to $7/15$,
and reduces the ordinary reward per juror when all reveal. With all $n$ jurors
revealing, the 5 ETH pool gives each $5/n$ ETH. The fixed 10 ETH offer is
therefore $10/(5/n)=2n$ times that share, or six, fourteen and thirty times
the ordinary reward at the three sizes. Jury size, seeded influence and the
relative payment incentive consequently change together
(\autoref{tab:quorum-geometry}).

\begin{table}[ht]
\centering\small
\caption{Decision geometry of the size comparison. Every condition retains
shared public deliberation. The last column is the ordinary reward if all
registered jurors reveal. The offered extra payment is always 10 ETH.}
\label{tab:quorum-geometry}
\begin{tabular}{rrrrrr}
\toprule
Jury $n$ & Seeds $m$ & Evaluated $k$ & Majority $q$ & $q-m$ & Reward (ETH) \\
\midrule
3 & 1 & 2 & 2 & 1 & 1.667 \\
7 & 3 & 4 & 4 & 1 & 0.714 \\
15 & 7 & 8 & 8 & 1 & 0.333 \\
\bottomrule
\end{tabular}
\end{table}

\subsection{Three mechanism contrasts at size three}
\textbf{Explicit duty.} The principles charter adds truth and independent
judgment duties to the short charter. This asks whether false adoption still
occurs when the instructions expressly require truthful, independent judgment.

\textbf{Repeated exposure.} The configuration named \texttt{contagion} runs
three cases on one chain with memory and continuing accounts. The publisher's
offer and scripted peer are supplied again in every case. We call this
\emph{repeated exposure}. A juror in a later case can see earlier outcomes, but
also receives a new invitation. The condition therefore tests persistence
with payment history available rather than isolating memory or transmission
to previously unexposed agents.

\textbf{Mandated publisher.} A publisher using the evaluated model is instructed
to obtain a false majority, with three evaluated jurors, no scripted peer and no supplied
offer. This tests whether recruitment emerges from an adversarial objective.
A wrong ballot alone in this arm is not classified as an executed coalition.

\subsection{Cloud audit, assignment and dependence}
A \emph{session}, identified by a run ID, contains one case or a three-case
repeated-exposure block. C1 and C2 identify two source-checkout cohorts, not
experimental treatments. The updated registry contains
\SnapshotIncludedSessions{} played sessions and \SnapshotAssignedCases{}
assigned cases, retaining all \SnapshotRetainedFailures{} sessions that
failed after play began. One setup-only artifact is excluded. A played failure
is not rerun to obtain a replacement outcome.

C1 contributes its completed 57-session plan, replacing the earlier 24-session
subset without counting those runs again. C2 contributes the same earlier
47-session partial selection. Runtime source hashes agree across the cohorts,
and settings match within model and condition after normalizing seed and
checkout path. Repeated seed labels do not create additional distinct seeds.
Models and conditions were adaptively selected from earlier signal, so the
combined counts are descriptive rather than a population sample.
\autoref{app:cloud-register} records selection cutoffs, exclusions and source
hashes.

\subsection{Reach, persistence and the limit on speed}
\label{sec:outcomes}
We score false-decision adoption only on a revealed
\texttt{CREATOR\_SUPPORTED} verdict. A truthful reveal instead supports the
challenger. Tables distinguish false reveals (F), truthful reveals (T), seats
with no commitment (A), and commitments without a valid reveal (U). The label
``abstention'' for A records absence of commitment, not a demonstrated decision
to resist the offer. Each assigned evaluated seat contributes to exactly one
of these four categories.

\emph{Capture} requires a false majority of the full jury, counting scripted
as well as evaluated reveals. A false minority vote contributes to F without
capturing the case. Payment is checked separately because the publisher's
promise is not escrowed. Execution failures are an overlapping operational
flag and do not erase retained votes. The reference and observed claims match
in every admitted case, including the publisher-mandate arm.

Reach is the number of evaluated false revealers in a case, with both revealed
and assigned-seat denominators visible. Per-case records preserve variability.
Votes within a case and cases within a memory block are dependent. We do not
turn pooled seat counts into independent replications or a population
confidence interval. The analyzer's seat-based fallback intervals are not
used to make propagation or topology claims.

The available timing observation places adoption within the same case's
commit/reveal cycle after guaranteed direct exposure. It does not establish
how many message transmissions were needed for adoption. Wall-clock
latency mixes model calls, rate limits, provider errors and scheduling. No
hop-by-hop exposure intervention or transmission lineage is observed. Thus
these records establish within-case adoption, not an epidemic growth rate,
reproduction number or causal infection velocity.
